\documentclass[11pt]{article}
\usepackage{fontspec}
\usepackage[french]{babel}
\usepackage{amsmath}
\usepackage{amssymb}
\usepackage{amsthm}
\usepackage{mathrsfs}
\usepackage{tikz-cd}
\usepackage{enumitem}
\usetikzlibrary{babel}

\newcommand{\printedpage}[1]{\marginpar{\footnotesize p.~#1}}
\newcommand{\msp}[1]{\marginpar{\footnotesize [#1]}}
\newcommand{\cercle}[1]{\tikz[baseline=(c.base)]\node[draw,circle,inner sep=1pt](c){#1};}
\newcommand{\incl}{\lhook\joinrel\longrightarrow}
\newcommand{\card}{\operatorname{card}}
\newcommand{\Ob}{\operatorname{Ob}}
\newcommand{\Ind}{\operatorname{Ind}}

\begin{document}

\printedpage{109}
% head: A. Grothendieck : Dérivateurs, Ch. XVIII, Édition des Universités de Montpellier II et Paris VII

D'autre part, chaque $\underline{S}_\alpha$ étant stable dans $\mathcal{M}$ par $\varinjlim$ grandes devant $\pi'_0$, il en est de même de leur intersection $\underline{S}$. A fortiori, celle-ci est stable par $\varinjlim$ grandes devant $\pi$, donc le foncteur pleinement fidèle canonique $\Ind_\pi(\underline{S}(\pi))\longrightarrow\mathcal{M}$ se factorise par $S$ en
\[ \Ind_\pi\underline{S}(\pi) \incl \underline{S} \; , \]
il faut prouver que ce foncteur est essentiellement surjectif. Soit donc $x\in S=\Ob\underline{S}$, tout revient à prouver que
\[ \varinjlim_{i\in\underline{S}(\pi)/x} x_i \xrightarrow{\ \sim\ } x \; , \tag*{$(*)$} \]
et que $\underline{S}(\pi)/x$ est grande devant $\pi$.

Notons que, comme $\pi'_0\geqslant\pi_0$, et $\pi^{\pi'_0}=\pi$, on aura $\pi^{\pi_0}=\pi$, donc $\pi$ est cardinal utile pour $\mathcal{M}$, donc $I=\mathcal{M}_\pi/x$ est grande devant $x$, et on a % [sic] grande devant $x$ -> grande devant $\pi$, apparently intended (one speaks of being large before a cardinal)
\[ \varinjlim_{i\in I=\mathcal{M}_\pi/x} x_i \xrightarrow{\ \sim\ } x \; . \tag*{$(**)$} \]
Donc, si on admet que $\underline{S}(\pi)/x$ est grande devant $\pi$,
\msp{page 294}
le fait que $(*)$ soit un isomorphisme équivaut à dire que le foncteur canonique
\[ J \overset{\text{déf}}{=} \underline{S}(\pi)/x \incl \mathcal{M}_\pi \; = \; I \] % [sic] $\mathcal{M}_\pi = I$ -> $\mathcal{M}_\pi/x = I$, since $I=\mathcal{M}_\pi/x$ above and in the display for $J_\alpha$ below
est \textit{cofinal} (cf. 4.4.2, p. 79). Mais on voit tout de suite que ce foncteur est un foncteur d'inclusion d'une sous-catégorie strictement pleine (stable par $\varinjlim$ grandes devant $\pi'_0$, de cardinal $\leqslant\pi$).

Notons maintenant que la flèche $(**)$ similaire à $(*)$, avec $\underline{S}(\pi)$ remplacée par $\underline{S}_\alpha(\pi)$, est bel et bien un isomorphisme, puisque $x\in\Ob\underline{S}\subset\Ob\underline{S}_\alpha$, et que $\underline{S}_\alpha$ est $\pi$-accessible. Cela s'interprète par le fait que le foncteur d'inclusion
\[ J_\alpha \overset{\text{déf}}{=} \underline{S}_\alpha(\pi)/x \incl \mathcal{M}_\pi/x \; = \; I \]
est cofinal. C'est encore une inclusion d'une sous-catégorie $J_\alpha$ strictement pleine de $I$, stable par $\varinjlim$ grandes devant $\pi'_0$ et de cardinal $\leqslant\pi$, et on a
\[ J \; = \; \bigcap J_\alpha \; . \]

Donc la condition voulue résultera du
\msp{page 295}

% marginal hand-drawn mark (manuscript) in the left margin beside the following lemma
\textbf{Lemme 4.13.10} (\footnote[125]{Il faut supposer que $\pi$ ne soit pas un cardinal limite, pour pouvoir conclure que $I$ est stable par $\varinjlim$ de type $I_{\alpha_\pi}$, et les $J_r$ stables dans $I$ par ces limites.},\footnote[126]{Cf. forme plus générale du lemme dans 4.13.10.2, p. 302.}). \textit{Soient $\pi'_0$ et $\pi>\pi'_0$ des cardinaux infinis, $I$ une catégorie équivalente à une petite catégorie, $(J_\alpha)_{\alpha\in A}$ une famille de sous-catégories $J_\alpha$ de $I$. On suppose :}
\begin{itemize}
\item[\cercle{a}] \textit{$I$ est grande devant $\pi$ (}\footnote[127]{Il suffit que $I$ soit grande devant tout $\pi'<\pi$.}\textit{).} % footnote 127 printed at the foot of p. 110
\item[\cercle{b}] \textit{$I$ est stable par $\varinjlim$ de cardinal $\leqslant\pi$, et grandes devant $\pi'_0$ (}\footnote[128]{Il suffit ici la stabilité par $\varinjlim$ de type $I_{\alpha_\pi}$, cf. page 302 [?].}\textit{).} % footnote 128 printed at the foot of p. 110
\end{itemize}
% foot: Version 17 décembre 2025

\printedpage{110}
% head: A. Grothendieck : Dérivateurs, Ch. XVIII, Édition des Universités de Montpellier II et Paris VII
\begin{itemize}
\item[\cercle{c}] \textit{$\forall\ r\in A$, $J_r$ est stable dans $I$ par les limites inductives envisagés dans b) (en particulier, $J_r$ est une sous-catégorie} strictement pleine \textit{de $I$).} % [sic] envisagés -> envisagées, agreement with « limites »
\item[\cercle{d}] \textit{$\forall\, r\in A$, $J_r$ est une sous-catégorie} cofinale \textit{de $I$ (}grande devant $\pi$\textit{).}
\item[\cercle{e}] $\card A\leqslant\pi$.
\end{itemize}
\textit{Soit $J=\bigcap_{r\in A}J_r$. Alors $J$ est encore cofinale dans $I$, et grande devant $\pi$ (et de plus stable dans $I$ par $\varinjlim$ de cardinal $\leqslant\pi$ et grandes devant $\pi'_0$).}

En d'autres termes, si $I$ est grande devant $\pi$, l'intersection $J$ d'une famille $(J_\alpha)_{\alpha\in A}$ de sous-catégories strictement pleines cofinales et grandes devant $\pi$, est encore
\msp{page 296}
cofinale et grande devant $\pi$, \textit{pourvu que} $\card A\leqslant\pi$, et que l'on postule une propriété de stabilité de $I$ et des $J_\alpha$ dans $I$, pour les limites inductives de cardinal $\leqslant\pi$ et grandes devant un cardinal convenable $\pi'_0<\pi$ donné.

\textbf{4.13.10.1.} Notons d'ailleurs que si $f : J\longrightarrow I$ est un foncteur cofinal \textit{pleinement fidèle}, avec $I$ grande devant $\pi$, alors $J$ est automatiquement grande devant $\pi$. Cela résulte p.ex. de 4.9.6.3 2° (p. 225) par le critère $\mathrm{c}_\pi$ (vu que pour $f$ pleinement fidèle, la condition F $3_\pi$ est tautologiquement satisfaite). On peut aussi le voir ainsi : Soit $B$ un ensemble avec $\card B\leqslant\pi$, il faut prouver que $J\longrightarrow J^B$ est cofinal, or on a
\[
\begin{tikzcd}
J \arrow[r, "\text{cof.\,?}"] \arrow[d, "\text{cof.}"'] & J^B \arrow[d, "\text{cof.}"] \\
I \arrow[r, "\text{cof.}"] & I^B \; ,
\end{tikzcd}
\]
où trois des quatre flèches du carré sont cofinales, $J^B\longrightarrow I^B$ l'étant comme produit de foncteurs cofinaux entre catégories \textit{filtrantes} (cf. 4.9.6.5, page 227) -- on sait en effet déjà que $J$ est filtrante.

\msp{page 297}
Donc le composé
\[ J \longrightarrow J^B \xrightarrow{\text{cof.}} I^B \]
est cofinal. Mais comme le deuxième foncteur $J^B\longrightarrow I^B$ est pleinement fidèle, il s'ensuit aussitôt que $J\longrightarrow J^B$ est cofinal (car $\xi\backslash J$ pour $\xi\in\Ob J^B$ est le même, à isomorphisme près, qu'on considère $\xi$ comme objet de $J^B$, ou de $I^B$).

Ainsi dans l'énoncé de la condition d) du lemme, il est inutile de spécifier que $J_\alpha$ est grande devant $\pi$, et de même dans la conclusion sur $J$.

Comme $J\incl I$ est pleinement fidèle, pour voir que ce foncteur est cofinal, il suffit de voir qu'il satisfait F 1, i.e. que pour tout $i\in\Ob I$, existe $j\in J$ avec $i\longrightarrow j$. On va construire $j$ comme limite d'un système inductif ordinal d'objets $i_\alpha$ de $I$, indexé par les ordinaux $\alpha$ tels que $\card\alpha<\pi$. Désignons par $\alpha_\pi$ le plus petite ordinal tel que % [sic] le plus petite ordinal -> le plus petit ordinal, agreement with « ordinal »
% foot: Version 17 décembre 2025
\printedpage{111}
% head: A. Grothendieck : Dérivateurs, Ch. XVIII, Édition des Universités de Montpellier II et Paris VII
\msp{page 298}
% marginal hand-drawn mark (manuscript) in the left margin beside the following two lines
$\card\alpha_\pi=\pi$. L'ensemble ordonné $I_{\alpha_\pi}$ des ordinaux $\alpha<\alpha_\pi$ étant grande devant $\pi'_0$ (\footnote[129]{faux si $\pi$ est un cardinal limite.}) et de cardinal $\pi$, on pourra donc prendre la limite $\varinjlim\limits_{\alpha\in I_{\alpha_0}} i_\alpha$ dans $I$. % [sic] $I_{\alpha_0}$ -> $I_{\alpha_\pi}$, the index set just defined is $I_{\alpha_\pi}$

Construction de $(i_\alpha)_{\alpha\in I_{\alpha_0}}$. Soit % [sic] $I_{\alpha_0}$ -> $I_{\alpha_\pi}$, the index set just defined is $I_{\alpha_\pi}$
\[ \alpha_\pi \; = \; \text{plus petit ordinal } \alpha \text{ tel que } \card\alpha=\pi, \tag*{(1)} \]
donc c'est l'ordinal caractérisé par
\[ \left\{
\begin{aligned}
&\card\alpha_\pi \; = \; \pi \\
&\alpha<\alpha_\pi \;\; \Longrightarrow \;\; \card\alpha \; < \; \pi \; ,
\end{aligned}
\right. \tag*{(2)} \]
% marginal hand-drawn mark (manuscript) in the left margin beside the following paragraph
ou encore par le fait que l'ensemble ordonné $I_{\alpha_\pi}$ des ordinaux $\alpha<\alpha_\pi$ est \textit{parfaitement bien ordonné}. Ainsi, $I_{\alpha_\pi}$ \textit{est grande devant $\pi'$ pour tout cardinal $\pi'<\pi$} (\footnote[130]{Faux si $\pi$ est cardinal limite.}). Dans les conditions b) c) de 4.3.10, il suffira de supposer $I$ stable par $\varinjlim$ de type $I_{\alpha_\pi}$, et les $J_r$ ($r\in A$) stables dans $I$ par ce type de limites. % [sic] 4.3.10 -> 4.13.10, the conditions b) c) are those of Lemme 4.13.10

\textbf{Lemme 4.13.10.1.} \textit{Soit $\pi$ un cardinal infini, $A$ un ensemble non vide tel que} $\card A\leqslant\pi$. \textit{Alors}
\msp{page 299}
\textit{il existe une application}
\[ \varphi \; : \; I_{\alpha_\pi} \longrightarrow A \]
\textit{telle que pour tout $r\in A$, $\varphi^{-1}(\{r\})$ soit une partie cofinale de $I_{\alpha_\pi}$, ou ce que est plus précis,} $\card\varphi^{-1}(\{r\})=\pi$. % [sic] ce que est -> ce qui est, relative pronoun subject

En effet, il suffit de considérer un ensemble $A'$ de cardinal $\pi$ (p.ex. $A'=I_{\alpha_\pi}$), l'application
\[ \psi \; = \; \mathrm{pr}_1 \; : \; A\times A'\longrightarrow A \; , \]
et de se donner une bijection
\[ A\times A' \xrightarrow{\ \sim\ } I_{\alpha_\pi} \]
laquelle existe puisque
\[ \card\,(A\times A') \; = \; \card A\cdot\card A' \; = \; \card A' \; = \; \pi \; = \; \card I_{\alpha_\pi} \; , \]
puisque $0\neq\card A\leqslant\card A'=\pi$ et $\pi$ infini. On suppose choisie $\varphi$ comme dans 4.13.10.1, et on construira le système inductif $(i_\alpha)_{\alpha\in I_{\alpha_\pi}}$ de telle façon qu'on ait
\[ i_\alpha \; \in \; J_{\varphi(\alpha)} \; . \]

Posant
\[ j \; = \; \varinjlim_{\alpha\in I_{\alpha_\pi}} i_\alpha \; , \]
% foot: Version 17 décembre 2025
\printedpage{112}
% head: A. Grothendieck : Dérivateurs, Ch. XVIII, Édition des Universités de Montpellier II et Paris VII
\msp{page 300}
pour tout $r$, comme
\[ I_r \overset{\text{déf}}{=} \varphi^{-1}(r) \; \subset \; I_{\alpha_\pi} \]
est cofinal dans $I_{\alpha_\pi}$, on aura donc
\[ j \; = \; \varinjlim_{\alpha\in I_r} i_\alpha \; , \]
et en vertu de $(*)$, les $i_\alpha$ pour $\alpha\in I_r$ sont dans $J_r$, donc par passage à la limite, utilisant l'hypothèse c), on trouve que $j\in J_r$, donc
\[ j \; \in \; \bigcap_{r\in A} J_r \; = \; J \; , \]
et comme on choisira
\[ i_0 \; = \; i \; , \tag*{$(**)$} \]
on aura bien trouvé $i\longrightarrow j$ avec $i\in J$. Donc il suffit de construire le système inductif $(i_\alpha)$ satisfaisant $(*)$ et $(**)$. % [sic] avec $i\in J$ -> avec $j\in J$, it is $j$ that was shown to lie in $J$

Si le système inductif est construit jusqu'au cran $i_\alpha\in\Ob J_{\varphi(\alpha)}$, on continue par
\[ i_\alpha \longrightarrow i_{\alpha+1} \; \in \; \Ob J_{\varphi(\alpha+1)} \; , \]
ce qui est possible puisque $J_{\varphi(\alpha+1)}$ est
\msp{page 301}
cofinale. Donc il reste à construire $i_\alpha$ pour $\alpha$ ordinal limite. Or l'ensemble des $i_\beta$ pour $\beta<\alpha$ est de cardinal $<\pi$, et comme par a) $I$ est grande devant $\pi$ (il suffit qu'il le soit devant tout cardinal $<\pi$), il existe un $i\in\Ob I$ qui majore tous les $i_\beta$. Choisissons des
\[ i_\beta \xrightarrow{v_\beta} i \; . \]

On n'aura pas a priori commutativité dans
\[
\begin{tikzcd}
i_{\beta'} \arrow[r, "\tau_{\beta,\beta'}"] \arrow[rr, bend right=15, "v_{\beta'}"'] & i_\beta \arrow[r, "v_\beta"] & i \; ,
\end{tikzcd}
\] % ?? uncertain: line style of the curved arrow -- solid or dotted
mais comme l'ensemble des double-flèches $i_{\beta'}\mathrel{\overset{v_\beta\tau_{\beta,\beta'}}{\underset{v_{\beta'}}{\rightrightarrows}}} i$ est de cardinal $<\pi$, $I$ étant grande devant $\pi$, il existe $i\longrightarrow i'$ qui les égalise tous. Comme $J_{\varphi(\alpha)}$ est cofinal dans $I$, il existe $i''\in J_{\varphi(\alpha)}$, et $i\longrightarrow i''$. On prend $i_\alpha=i''$, avec les homomorphismes composés
\[
\begin{tikzcd}
i_\beta \arrow[r, "v_\beta"] \arrow[rrr, bend right=12, "\tau_{\alpha,\beta}"'] & i \arrow[r] & i' \arrow[r] & i''
\end{tikzcd}
\] % ?? uncertain: line style of the curved arrow -- solid or dotted
% foot: Version 17 décembre 2025
\printedpage{113}
% head: A. Grothendieck : Dérivateurs, Ch. XVIII, Édition des Universités de Montpellier II et Paris VII
\msp{page 302}
comme morphismes de transition. Cela achève la démonstration de 4.13.10 et par là de 4.13.9, 4.13.8.

\textbf{Remarque 4.13.10.2.} Voici ce qu'on a eu à utiliser vraiment, pour établir la conclusion de 4.13.10, en termes de $I$, $\pi$, $(J_r)_{r\in A}$ (et sans introduire de $\pi'_0$).
\begin{itemize}
\item[\cercle{a$'$}] $I$ est grande devant tout cardinal $\pi'<\pi$.
\item[\cercle{b$'$}] $I$ est stable par $\varinjlim$ de type $I_{\alpha_\pi}$.
\item[\cercle{c$'$}] $\forall\, r\in A$, $J_r$ est stable dans $I$ par les $\varinjlim$ de type $I_{\alpha_\pi}$.
\item[\cercle{d$'$}] $\forall\, r\in A$, $J_r$ est cofinale dans $I$.
\item[\cercle{e$'$}] $\card A\leqslant\pi$.
\end{itemize}
(\textbf{N.B.} d$'$) et e$'$) ne sont autres que d) et e) de 4.3.10, tandis que a$'$), b$'$), c$'$) sont des versions affaiblies de a), b), c) respectivement.) % [sic] 4.3.10 -> 4.13.10, the conditions d) e) are those of Lemme 4.13.10

\textbf{4.13.8 (reformulation).} \textit{Soit pour $r\in A$, $\pi''_r\geqslant\pi'_r\geqslant\pi_0$ (= cardinal utile pour $\mathcal{M}$ donné) comme spécifié dans 4.13.2 b) (et c), soit $\pi'=\sup_{r\in A}\pi'_r$, $\pi''=\sup_{r\in A}\pi''_r$, de sorte que pour tout cardinal $\pi$ avec $\pi^{\pi'_0}=\pi\geqslant\pi''$, $\underline{S}_r$ soit $\pi$-accessible dans $\mathcal{M}$ (et $\pi$ utile pour $\underline{S}_r$} et \textit{pour $\mathcal{M}$). Soit $c$ un cardinal} non limite \textit{tel que $c\geqslant$} $\card A$ \textit{et $c>\pi'$ (p.ex. le successeur de} $\sup(\card A,\pi')$\textit{). Prenons $\pi$ tel que}
\[ \pi^{\pi'} \; = \; \pi \; \geqslant \; \sup(\pi'',c) \; . \]
\textit{Alors $\underline{S}=\bigcap_{r\in A}\underline{S}_r$ est $\pi$-accessible.}

En effet, si $\alpha_c$ est le premier ordinal de cardinal $c$, comme $c$ n'est pas limite, $I_{\alpha_c}$ est grande devant tout cardinal $c'<c$, en particulier devant $\pi'$. Or pour tout $x\in\Ob\mathcal{M}$, $I=\mathcal{M}_\pi/x$ et les $J_r=\underline{S}_r/x$ sont stables par $\varinjlim$ grandes devant $\pi'$ et de cardinal $\leqslant\pi$,
\msp{page 302 bis}
et les $J_r$ stables dans $I$ pour ces limites, donc il en est en particulier ainsi pour les $\varinjlim$ de type $I_{\alpha_c}$. On peut donc appliquer le résultat 4.13.10.2 pour conclure que si les $J_r$ sont cofinaux dans $I$ (qui est grande devant $\pi$, $\pi$ étant utile pour $\mathcal{M}$), alors $J=\underline{S}/x$ l'est également. Cela implique alors, comme on a vu, que $\underline{S}$ est $\pi$-accessible dans $\mathcal{M}$.

\textbf{4.13.10.3.} On peut reformuler la reformulation précédente, en utilisant la terminologie introduite plus bas (4.13.23) :

\textit{Soit $c$ un cardinal, $\alpha_c$ le plus petit ordinal $\alpha$ tel que} $\card\alpha=c$, \textit{$\pi$ un cardinal $\geqslant c$. On suppose}
\begin{itemize}
\item[$(*)$] \textit{$\pi$ est utile pour $\mathcal{M}$, et $\mathcal{M}$ est stable par $\varinjlim$ de type $I_{\alpha_c}$. (Cette deuxième condition est le cas p.ex. si $c$ est} non limite \textit{et $c>\pi_0$, $\pi_0$ étant un cardinal tel que $\mathcal{M}$ satisfait \textit{2}$_{\pi_0}$, p.ex. un cardinal} utile \textit{pour $\mathcal{M}$ (}\footnote[131]{Car alors $I_{\alpha_c}$ est grande devant $\pi_0$.}\textit{)).}
\end{itemize}
\textit{Alors l'ensemble des sous-catégories $\underline{S}$ de $\mathcal{M}$ qui sont $\alpha_c$-$\pi$-accessible dans $\mathcal{M}$ (i.e. qui sont $\pi$-accessibles et stables dans $\mathcal{M}$ par $\varinjlim$ de type $I_{\alpha_c}$) est stable par intersection de cardinal $\leqslant c$. Plus généralement, si $\Phi$ est un ensemble de catégories telles que} % [sic] accessible -> accessibles, agreement with « sous-catégories »
\begin{itemize}
\item[\textit{1°)}] \textit{$I\in\Phi\Longrightarrow$} $\card I\leqslant\pi$ \textit{et}
\item[\textit{2°)}] \textit{$I_{\alpha_c}\in\Phi$,}
\end{itemize}
\textit{alors l'ensemble des sous-catégories de $\mathcal{M}$ qui sont $\Phi$-$\pi$-accessibles dans $\mathcal{M}$, est stable par intersection de cardinal $\leqslant c$.}
% foot: Version 17 décembre 2025

\end{document}
